\subsection{THETA3-A16: the literal full host has an abelian-field restriction}\label{A16}

The distinction at the end of A11 is structural. An arbitrary
CM field cannot be enlarged into the field of the literal
\texttt{Formalism.BasicNotions.Host}. This rules out one
overly broad matching procedure; it does not contradict the
general oscillator identities or the conditional constructions
A13--A15 on an existing full host.

\paragraph{The exact field restriction in the source.}
The pinned archive defines \texttt{HostRaw} in
\texttt{Formalism/SharpRestriction/RamificationFrobenius/}
\texttt{TotientGaloisClasses/ConjEMemHostFIff.lean},
lines 12--47. Its fields include
\[
 \operatorname{Gal}(E/\Q)\simeq\phi.\mathrm{GG},
 \qquad
 E\hookrightarrow\Q(\zeta_{\mathrm{QS}\,q\,d}),
 \tag{A16a}
\]
as well as the explicit assertion that $E/\Q$ is Galois.
The group $\phi.\mathrm{GG}$ is
$\{\pm1\}^2\times\phi.\mathrm{MA}$, where
$\phi.\mathrm{MA}$ is the multiplicative version of a finite
abelian additive group. The exact definitions are
\texttt{ZsmulGeneratorEqZeroOfProjection.lean},
lines 15--23, 29--33, 87--89, in the same directory;
\texttt{Formalism/FrobeniusOrbits/Statement.lean},
line 27, defines the two sign factors;
\texttt{FrameFacts.lean}, lines 30--35, also expands the
two sign factors directly. Hence $E/\Q$ is abelian.
This conclusion follows directly from the displayed group
isomorphism; it does not require any converse to the
Kronecker--Weber theorem.

The record \texttt{HostAx} extends \texttt{HostRaw} in
\texttt{ConjEMemHostFIff.lean}, line 174. The stronger
\texttt{Host} extends \texttt{HostAx} in
\texttt{Formalism/LocalRingCardinality/DataHost.lean},
line 39. Thus every full host used by the stated
\texttt{D7}/\texttt{StdPair} construction has the restriction
(A16a). Enlarging the separate coefficient field
\texttt{LL}, as in A13--A14, does not change this host field.

\paragraph{A concrete CM field excluded by every such host.}
Put
\[
 F=\Q(\sqrt2),\qquad
 K=F(\alpha),\qquad
 \alpha^2=-(3+\sqrt2).
 \tag{A16b}
\]
Then $K$ is a quartic CM field which is not Galois over
$\Q$. Indeed, both real embeddings of $F$ send
$3+\sqrt2$ to a positive number. Its negative is not
a square in $F$, since it is negative at either real
embedding. Consequently $[K:F]=2$, and $K/F$ is a
totally imaginary quadratic extension of the totally
real field $F$. Also $\sqrt2=-\alpha^2-3$, so
$K=\Q(\alpha)$ and its minimal polynomial is
\[
                     X^4+6X^2+7.
 \tag{A16c}
\]

Fix the embedding with $\alpha=i\sqrt{3+\sqrt2}$.
Every element of $K$ is uniquely $a+b\alpha$ with
$a,b\in F\subset\R$, and it is real exactly when
$b=0$. Thus $K\cap\R=F$. If $K/\Q$ were normal,
it would contain another root $\beta$ of (A16c) with
$\beta^2=-(3-\sqrt2)$. Then
$(\alpha\beta)^2=7$, so $\sqrt7\in K\cap\R=F$.
This is impossible: if $\sqrt7=a+b\sqrt2$ with
$a,b\in\Q$, squaring gives $2ab=0$, and then either
$a^2=7$ or $b^2=7/2$, neither of which has a rational
solution. This proves the claimed nonnormality.

No $\Q$-embedding of $K$ into a full-host field exists.
Every subfield of a finite abelian Galois extension is
Galois over $\Q$: its corresponding subgroup is normal,
and the Galois correspondence applies. An embedded copy
of $K$ would therefore be Galois, contrary to the proof
above. More strongly, for every finite extension
$K'/K$ there is no $\Q$-embedding $K'\hookrightarrow E$
into a full-host field, since its restriction would
embed $K$.

The same obstruction applies if one starts with a
Galois CM field as in 032. For completeness, the normal
closure here is
\[
       \widetilde K=K(\sqrt7),\qquad
       \widetilde F=\Q(\sqrt2,\sqrt7).
 \tag{A16d}
\]
The field $\widetilde F$ is totally real of degree four;
$\widetilde K=\widetilde F(\alpha)$ is its totally
imaginary quadratic extension. The four roots
$\pm\alpha,\ \pm\sqrt7/\alpha$ of (A16c) all lie
in it, so it is the normal closure and is a Galois
CM field of degree eight. It is nonabelian, since it
contains the nonnormal subfield $K$. It cannot embed
in a full-host field, nor can any finite extension
containing it. The four-factor theorem in 032,
Theorem 2.5 and its reduction Proposition 2.6, printed p.8,
is stated for every Galois CM field; it is not
restricted to the abelian fields of (A16a).

\paragraph{What is refuted, and the remaining comparison.}
The following proposed universal step is refuted:
start with an arbitrary 032 CM field, retain its field
action by a field embedding, enlarge it finitely, and
obtain the literal full-host record of the recovered
arithmetic archive. The explicit counterexample is
(A16b), or the Galois example (A16d). This is stronger
than an unresolved search for auxiliary primes or
for a factorization of a different.

It does not refute A11's generic
\texttt{Geo.Tor.Host}, which has different fields and
hypotheses. It does not refute A13--A15, which start
with an existing full host, or B25's finite definite
period identity on the actual oscillator spaces.
It does not assert the impossibility of an unrelated
correspondence, a different field action, or a
generalized arithmetic theorem. The viable general
route is to use the field-independent proved
oscillator and finite-transform identities with
their own explicit hypotheses, and prove the needed
arithmetic estimates for the given CM field. The
literal archive theorem remains usable on its actual
abelian-field domain.

In the requested Grothendieck framework, an
$E$-linear correspondence must retain its specified
endomorphism-field action on the realization. A
change of coefficient field is distinct from a
change of that action. This is the realization and
correspondence setting of \emph{cote44}, scans 7--8;
the algebraic-projector discussion of \emph{cote16},
scans 43--46, supplies no embedding of a nonnormal
CM field into an abelian extension.
